\documentclass{article}
\usepackage{graphicx}
\usepackage[margin=1.5cm]{geometry}
\usepackage{amsmath}

\begin{document}
\twocolumn

\title{Wednesday Warm Up: Unit 0, Electrostatics I and II}
\author{Prof. Jordan C. Hanson}

\maketitle

\section{Memory Bank}

\begin{itemize}
\item $\vec{F} = q \vec{E}$ ... Force on a charge $q$ in the presence of an $\vec{E}$-field.
\item $\vec{F} = m\vec{a}$ ... Newton's 2nd Law.
\item $\vec{E}_{\rm plane} = \frac{\sigma}{2\epsilon_0}\hat{k} = \frac{Q}{2A\epsilon_0}\hat{k}$ ... The electric field of a plate or disk with charge $Q$ and area $A$, for distances much smaller than the radius or length.
\item \textbf{Mass of electron}: $9.11\times 10^{-31}$ kg.
\item \textbf{Charge of electron}: $1.6\times 10^{-19}$ C.
\item $\Delta U = q\Delta V$ ... Relationship between change in potential energy, charge, and voltage.
\item $V(r) = k q / r$ ... Potential (voltage) of a point charge.
\item $V(x) = E x + V_0$ ... Linear relationship between voltage, electric field, and distance between two charge plates to form a capacitor.
\item $(1/2) m \bar{v^2} = (3/2) k T$.  Relationship between temperature and average atomic/molecular kinetic energy.  $T$ is the temperature in Kelvin, and $k = 1.38 \times 10^{-23}$ J/K, known as Boltzmann's constant.
\item $Q = C\Delta V$ ... Relationship between stored charge, capacitance, and voltage.
\end{itemize}

\section{Electric Fields}

\begin{enumerate}
\item An evacuated tube uses an accelerating voltage of 40 kV to accelerate electrons to hit a copper plate and produce x rays. Non-relativistically, what would be the maximum speed of these electrons? \\ \vspace{4cm}
\item Imagine a capacitor that has the following properties: the distance between the plates is 1 mm, the voltage in the middle is 0 Volts, $V_{\rm max} = 6$ Volts, $V_{\rm min} = -6$ Volts, and $\Delta V = 12$ V.  (a) What is the electric field in between the plates? (b) If an electron is accelerated by the capacitor voltage, what is the kinetic energy in eV? \\ \vspace{2cm}
\item Fusion probability is greatly enhanced when appropriate nuclei are brought close together, but mutual Coulomb repulsion must be overcome. This can be done using the kinetic energy of high-temperature gas ions or by accelerating the nuclei toward one another. (a) Calculate the potential energy of two singly charged nuclei separated by $1.0 \times 10^{-12}$ m by finding the voltage of one at that distance and multiplying by the charge of the other. (b) At what temperature will atoms of a gas have an average kinetic energy equal to this needed electrical potential energy? \\ \vspace{4cm}
\item Show that units of V/m and N/C for electric field strength are indeed equivalent. \\ \vspace{2cm}
\item What charge is stored in a 180 µF capacitor when 120 V is applied to it? \\ \vspace{2cm}
\end{enumerate}

\end{document}
